% Collocamento: il riepilogo del playtest di novembre 2026 per i game designer.
% Le immagini si preparano con py -3 playtest/figure.py; poi, da dentro playtest/,
% pdflatex feedback_playtest_2026-11.tex (due volte, per i riferimenti alle figure).
\documentclass[11pt,a4paper]{article}
\usepackage[utf8]{inputenc}
\usepackage[T1]{fontenc}
\usepackage[italian]{babel}
\usepackage{lmodern}
\usepackage[margin=2.5cm]{geometry}
\usepackage{graphicx}
\usepackage{booktabs}
\usepackage{tabularx}
\usepackage{xltabular}
\usepackage{array}
\usepackage{enumitem}
\usepackage[hidelinks]{hyperref}

\newcolumntype{L}{>{\raggedright\arraybackslash}X}
\newcolumntype{P}[1]{>{\raggedright\arraybackslash}p{#1}}
\setlength{\parskip}{0.5em}
\setlength{\parindent}{0pt}

\title{Collocamento: che cosa cambia dopo il playtest}
\date{5 novembre 2026}

\begin{document}
\maketitle

La partita si accorcia a tre set e perde i token obiettivo, il limite di mano e lo scorrimento del mercato.
Entrano un jolly, cioè l'Ambito giocato coperto, sei poteri pensati per l'interazione e, da provare, un mazzo
di eventi sul mercato del lavoro in Italia. Le modifiche sono già nei file del print and play; i PDF da stampare
si rifanno quando sarà decisa la versione nuova. In fondo ci sono le domande ancora aperte.

\section{Regole}

Il confronto è con il regolamento stampato il 28 ottobre.

{\small
\begin{tabularx}{\textwidth}{@{}P{2.8cm}LL@{}}
\toprule
Punto & Regolamento stampato & Dopo il playtest \\
\midrule
Fine della partita & Al primo quinto set si finisce il giro, poi si gioca un giro in più. & Quando qualcuno manda al lavoro la terza Persona, ogni altro partecipante gioca ancora un turno. Poi si contano i punti. \\
Token obiettivo & Quattro token: tre set della stessa categoria, quattro set di categorie diverse. & Tolti. Le categorie restano sulle carte, ma non hanno più effetto sul gioco. \\
Persone in tavola & Una per partecipante. & Invariato: il playtest lo conferma. \\
Mercato & Chi pesca fa uscire la Persona in fila da più tempo, e ne entra una nuova. & Non scorre: cambia solo quando qualcuno prende una Persona. \\
Limite di mano & 6 carte. & Nessun limite. \\
Chi inizia & Chi è più giovane. & Chi ha fatto per ultimo un colloquio di lavoro. \\
Set completo & Le Formazioni vanno negli scarti, Persona e Lavoro restano davanti a sé. & In più il set si chiude: il Lavoro va dritto sopra la Persona e le due carte si girano in orizzontale. \\
Ambito coperto & Non previsto. & Vale come un ambito a scelta, come Formazione o come Lavoro. Si dice quale quando lo si gioca. \\
Annuncio & Non previsto. & Ogni azione si annuncia prima di farla: serve al potere Contrordine. \\
\bottomrule
\end{tabularx}
}

\section{Carte e libretto}

\paragraph{Icone in alto, sempre nello stesso ordine.} Le tre icone stavano in colonna sul lato destro e, con le
carte a ventaglio in mano, la carta accanto le copriva. Ora stanno in fila in alto a sinistra (figura~\ref{fig:carte}),
sempre nello stesso ordine: Manualità, Digitale, Comunicazione, Organizzazione, Creatività, Ricerca. Così due carte si
confrontano a colpo d'occhio. Nelle pile le carte si sfalsano in verticale: ognuna sporge in alto da quella che ha
sopra, e resta visibile la fila di icone di tutte.

\begin{figure}[h]
\centering
\includegraphics[width=0.78\textwidth]{img/carte.jpg}
\caption{Due Persone e un Ambito con le icone in alto, nell'ordine fisso.}
\label{fig:carte}
\end{figure}

\paragraph{Un dorso nuovo per gli Ambiti: la Ruota.} Una carta Ambito coperta è un ambito qualsiasi, e il dorso
lo racconta: le sei icone delle competenze stanno su una ruota della fortuna, nello stesso ordine delle carte, con
la freccia che gira (figura~\ref{fig:ruota}). Dove si ferma, quello è l'ambito.

\begin{figure}[h]
\centering
\includegraphics[width=0.3\textwidth]{img/ruota.jpg}
\caption{Il dorso degli Ambiti.}
\label{fig:ruota}
\end{figure}

\paragraph{Libretto.} Il regolamento non spiega più come stampare e tagliare il print and play. Senza i token,
la pagina dei punti ospita anche la vittoria, e quella della fine partita un esempio di conto finale.

\section{Poteri}

Restano Passaparola e Corso intensivo, escono Ho un amico e Chiamata diretta, ne entrano quattro nuovi per dare
più interazione e spargere i punti. Con sei poteri, ognuno va su 10 Persone.

{\small
\begin{tabularx}{\textwidth}{@{}P{3.4cm}Lp{1.4cm}@{}}
\toprule
Potere & Che cosa fa & Stato \\
\midrule
Passaparola & Pesca 1 carta Ambito. & Resta \\
Corso intensivo & Gioca subito fino a 2 carte Ambito come Formazioni. & Resta \\
Networking & Il set di questa Persona vale 2 punti in più, contati a fine partita. & Nuovo \\
Concorrenza & Prendi una Formazione da una pila aperta di un altro partecipante, purché gliene resti almeno una, e mettila in una tua pila. & Nuovo \\
Cacciatore di teste & Prendi una Persona da una pila aperta di un altro partecipante e mettila su una tua pila libera che copre le sue ambizioni. Il suo potere non si attiva. & Nuovo \\
Contrordine & Fino al tuo prossimo turno, una volta: quando qualcuno annuncia la sua azione, puoi ordinargli di cambiarla. & Nuovo \\
Ho un amico & Prenota una Persona del mercato. & Esce \\
Chiamata diretta & Piazza un'altra Persona ignorando un'ambizione. & Esce \\
\bottomrule
\end{tabularx}
}

Una pila è aperta finché non ha un Lavoro. Per ora i poteri vanno alle Persone con lo stesso criterio di prima:
il più utile alle Persone più difficili da formare. Cacciatore di teste e Contrordine sono nomi proposti.

\section{Carte evento, da provare}

Gli eventi sono l'idea del playtest da provare per prima. Sono un mazzo di 12 carte. All'inizio di ogni giro chi ha
iniziato la partita scopre un evento, che vale fino al giro successivo. Ogni carta ha un effetto e una riga di spunto
su come funziona il mercato del lavoro in Italia: il gioco diventa un'occasione per parlarne. I due poteri tolti
tornano come eventi: Raccomandazione e Tirocinio.

{\small
\begin{xltabular}{\textwidth}{@{}P{2.9cm}LL@{}}
\toprule
Evento & Effetto per il giro & Spunto \\
\midrule
\endhead
PNRR & Chi apre o migliora una formazione può giocare un secondo Ambito come Formazione. & Con i fondi del PNRR il programma GOL finanzia corsi per chi cerca lavoro. \\
Supporto formazione e lavoro & Chi ha meno set chiusi pesca 1 Ambito all'inizio del suo turno. & Dal 2023 chi è in difficoltà economica e segue un corso può ricevere un'indennità mensile. \\
Raccomandazione & Chi inserisce una Persona può prenderne un'altra dal mercato e tenerla da parte, scoperta, fino al suo prossimo turno. Se non la inserisce, torna in fondo al mazzo. & Molte assunzioni passano da parenti, amici e conoscenti più che dai canali ufficiali. \\
Bonus assunzioni & Il primo set chiuso nel giro si prende questa carta: vale 2 punti a fine partita. & Gli incentivi per chi assume giovani, donne o al Sud sono tanti e cambiano quasi ogni anno. \\
Click day & Si rivelano 2 Persone in più. La prima voce che urla «Click!» ne mette una su una sua pila che la copre. A fine giro quelle rimaste tornano in fondo al mazzo. & Con il decreto flussi i permessi per lavorare in Italia vanno a chi invia per primo la domanda, nel giorno del click. \\
Tirocinio & Si può inserire una Persona ignorando una delle sue ambizioni; il suo potere non si attiva. & Molti tirocini pagano un rimborso basso e a volte prendono il posto di un contratto. \\
Centri per l'impiego & Il mercato non si ripristina: le Persone prese si sostituiscono tutte a fine giro. & I centri per l'impiego trovano lavoro a una piccola parte di chi lo cerca. \\
Fuga dei cervelli & Le Persone del mercato con l'icona Ricerca vanno in fondo al mazzo; se ne rivelano altre al loro posto. & Ogni anno migliaia di giovani laureati vanno a lavorare all'estero. \\
Sciopero generale & Nessuno chiude set. Chi pesca prende 3 Ambiti invece di 2. & Lo sciopero è un diritto scritto nella Costituzione, all'articolo 40. \\
Contratti a termine & I set chiusi nel giro restano dritti e aperti fino al giro dopo, Formazioni comprese. Se Cacciatore di teste ne porta via la Persona, il Lavoro torna in mano a chi l'aveva giocato. & La maggior parte dei nuovi contratti di lavoro è a tempo determinato. \\
Salario minimo & Il primo set da 5 punti chiuso nel giro, senza ambizioni che combaciano, si prende questa carta: vale 2 punti. & In Italia non c'è un salario minimo fissato per legge: i minimi li stabiliscono i contratti collettivi. \\
Mismatch & Non si può inserire una Persona su una pila con una sola Formazione. & Le imprese dicono di non trovare le persone giuste per molte delle assunzioni che prevedono. \\
\bottomrule
\end{xltabular}
}

Sette eventi aiutano, quattro frenano e uno, Contratti a termine, apre al furto. Gli spunti vanno controllati
sulle fonti prima di stampare le carte. L'altra idea del playtest, Formazioni da due competenze con i Lavori in un
mazzo separato, è stata scartata: cambiava troppo il gioco.

\section{Da decidere insieme}

\begin{enumerate}[leftmargin=*,itemsep=0.4em]
\item \textbf{Il jolly coperto va limitato.} 51 Persone su 60 hanno le stesse tre icone di un ambito. Giocato
coperto come Lavoro, quell'ambito chiude il set a 16 punti, perché la pila copre già quelle icone. Proposte: come
Lavoro vale 5 punti senza bonus, oppure si gioca coperto solo come Formazione.
\item \textbf{Il mercato può bloccarsi.} Lo scorrimento era stato aggiunto perché il mercato restava fermo su
Persone difficili: nelle simulazioni in tre dimezzava l'attesa, da 3,4 a 1,7 giri. Ora le Persone si possono
rubare; da guardare al prossimo playtest.
\item \textbf{Le categorie.} Senza token le cinque categorie degli Ambiti non hanno più effetto. Si tengono come
ambientazione o si usano altrove, per esempio negli eventi o nei poteri?
\item \textbf{Set chiusi.} Le Formazioni vanno ancora negli scarti, o restano nel set girato in orizzontale?
\item \textbf{Contrordine.} Chi lo subisce deve fare un'altra delle cinque azioni, o può rifare la stessa con
carte diverse?
\item \textbf{Concorrenza.} Può togliere una Formazione da una pila che ha già la Persona? Nella proposta sì: la
Persona resta, ma chiudere il Lavoro diventa più difficile.
\item \textbf{Eventi.} Valgono per un giro intero, come nella proposta, o solo per il turno di chi li scopre?
\item \textbf{Nomi.} Cacciatore di teste e Contrordine sono proposte: vanno scelti i nomi definitivi.
\end{enumerate}

\end{document}
